\documentclass[11pt,a4paper]{extarticle}
\usepackage[left=25mm,right=20mm,top=20mm,bottom=20mm]{geometry}
\usepackage{fontspec}
\IfFontExistsTF{Times New Roman}{\setmainfont{Times New Roman}}{\setmainfont{Liberation Serif}}

\usepackage{polyglossia}
\setdefaultlanguage{russian}
\setotherlanguage{english}
\usepackage{amsmath,amssymb,graphicx,booktabs,tabularx,array}
\usepackage{setspace,caption,float,enumitem,titlesec}
\usepackage[unicode,hidelinks]{hyperref}
\hypersetup{pdfauthor={Шешуков Евгений Игоревич},pdflang={ru-RU}}
\setstretch{1.02}
\setlength{\parindent}{1.25cm}
\setlength{\parskip}{3pt}
\setlength{\emergencystretch}{2em}
\setlength{\tabcolsep}{5pt}
\setlength{\intextsep}{8pt}
\renewcommand{\arraystretch}{1.18}
\captionsetup{font=small,labelsep=colon,justification=centering}
\setlist{nosep,topsep=5pt,leftmargin=1.4cm}
\titleformat{\section}{\large\bfseries}{\thesection.}{0.5em}{}
\titleformat{\subsection}{\normalsize\bfseries}{\thesubsection.}{0.5em}{}
\titlespacing*{\section}{0pt}{12pt}{7pt}
\titlespacing*{\subsection}{0pt}{9pt}{4pt}
\setlength{\abovedisplayskip}{7pt}
\setlength{\belowdisplayskip}{7pt}
\widowpenalty=10000
\clubpenalty=10000
\raggedbottom
\newcommand{\code}[1]{\texttt{\small #1}}
\begin{document}
\begin{center}{\Large\bfseries Компьютерная графика. Подготовка к защите}\\Лабораторные работы 3 и 4 · Вариант 13\end{center}
\section*{Лабораторная 3. Как превратить линию в пиксели}
Растеризация выбирает пиксели дискретной сетки, приближающие непрерывную фигуру. В программе точка $(x,y)$ задаёт центр пикселя. Для отрезка обозначим $\Delta x=x_1-x_0$, $\Delta y=y_1-y_0$ и $L=\max(|\Delta x|,|\Delta y|)$. Ведущей называют ось с большим модулем приращения. Шаг на единицу вдоль неё позволяет построить связный отрезок из $L+1$ пикселей, включая оба конца. При $L=0$ выводится одна точка.
\subsection*{Три способа построить отрезок}
\textbf{Пошаговый метод.} Если ведущая ось $X$, перебираем целые $x$ от начала к концу и каждый раз вычисляем $y=y_0+(x-x_0)\Delta y/\Delta x$. Ординату округляем. Для крутого или вертикального отрезка ведущей становится $Y$, и по ней вычисляем $x$.

\textbf{ЦДА (цифровой дифференциальный анализатор).} Начинаем в $(x_0,y_0)$. На каждом шаге прибавляем к координатам $(\Delta x/L,\Delta y/L)$. Округляем только при выводе пикселя. Если округлять сами накопленные координаты, дробная часть шага потеряется. Для $(0,0)\to(13,8)$ приращение равно $(1,8/13)$.

\textbf{Брезенхем.} Выбираем между двумя соседними пикселями по знаку целочисленной ошибки. Пусть $a$ и $b$ --- модули приращений ведущей и второй оси. Сначала $d=2b-a$. При $d<0$ шагаем только по ведущей оси и прибавляем к $d$ величину $2b$. При $d\ge0$ шагаем по обеим осям и прибавляем $2(b-a)$. Направления задаются знаками исходных приращений.

Для $(0,0)\to(13,8)$ начальная ошибка $d=3$. Первый шаг диагональный: получаем $(1,1)$ и $d=3-10=-7$. Второй шаг прямой: получаем $(2,1)$ и $d=-7+16=9$. Так работают все восемь октантов после выбора ведущей оси и знаков.
\subsection*{Окружность Брезенхема и сглаживание Ву}
Для окружности достаточно построить один октант. Стартуем с $(u,v)=(0,R)$ и $d=3-2R$. Пока $u\le v$, отражаем точку в $(\pm u,\pm v)$ и $(\pm v,\pm u)$ относительно центра. При $d<0$ прибавляем $4u+6$. Иначе прибавляем $4(u-v)+10$ и уменьшаем $v$ на 1. Затем увеличиваем $u$ на 1. В формулах используются старые координаты. На осях и диагоналях отражения совпадают, поэтому повторы убираются. Для $R=8$ получается 44 пикселя.

В схеме Ву точная ордината $y=j+f$, где $j=\lfloor y\rfloor$. Два соседних пикселя получают веса $1-f$ и $f$. Например, при $y=2.25$ пиксель на строке 2 получает вес 0.75, на строке 3 --- 0.25. Это смягчает ступеньки. Для крутой линии меняются роли осей.

\textbf{Что сказать о скорости.} Построение отрезка требует $O(L)$ шагов, окружности --- $O(R)$. Брезенхем использует целые сложения и сравнения. В лабораторной замеряется создание списка пикселей без рисования и записи шагов, берётся медиана пяти серий.

\clearpage
\section*{Лабораторная 3. Разбор вычислений для защиты}
\subsection*{Почему у трёх методов разные формулы?}
Для отрезка $(0,0)\to(13,8)$ идеальная ордината при $x=4$ равна $32/13\approx2.46$. Пошаговый метод получает её подстановкой $x$ в уравнение прямой. ЦДА приходит к тому же значению четырьмя прибавлениями $8/13$. Оба выбирают пиксель $(4,2)$. Брезенхем приходит в этот пиксель последовательностью решений по знаку ошибки.
\begin{center}\begin{tabular}{clrl}\toprule
Шаг & Текущий пиксель & $d$ & Следующий пиксель\\\midrule
0 & $(0,0)$ & 3 & $(1,1)$, диагональный шаг\\
1 & $(1,1)$ & $-7$ & $(2,1)$, прямой шаг\\
2 & $(2,1)$ & 9 & $(3,2)$, диагональный шаг\\
3 & $(3,2)$ & $-1$ & $(4,2)$, прямой шаг\\\bottomrule\end{tabular}\end{center}
\textbf{Откуда берётся начальное $d=2b-a$?} При первом шаге по ведущей оси идеальная линия поднимается на $b/a$. Чтобы выбрать между текущей строкой и следующей, сравниваем подъём с половиной пикселя. Умножение разности $b/a-1/2$ на $2a$ даёт $2b-a$. При следующем шаге к ошибке прибавляется $2b$. Если алгоритм поднялся на целую строку, дополнительно вычитается $2a$.

\textbf{Зачем выбирать ведущую ось?} Для отрезка $(0,0)\to(2,7)$ перебор только по $x$ дал бы три точки на восемь строк. Поэтому перебираем $y=0,1,\ldots,7$, а $x$ определяем по наклону. При движении влево или вниз шаг соответствующей координаты отрицательный.

\subsection*{Как вручную начать окружность}
При $R=8$ начальные значения $(u,v,d)=(0,8,-13)$. Ошибка отрицательна, значит сохраняем $v=8$, прибавляем к $d$ величину $4\cdot0+6=6$ и увеличиваем $u$. Получаем $(1,8,-7)$. Следующий шаг даёт $(2,8,3)$. Теперь ошибка положительна, поэтому уменьшаем $v$ и прибавляем $4(2-8)+10=-14$. Новое состояние $(3,7,-11)$.

Из точки $(3,7)$ получаем восемь отражений, включая $(7,3)$ и $(-3,7)$. Если центр окружности $(c_x,c_y)$, к каждой полученной паре прибавляем координаты центра. Для $(0,8)$ отражения дают только четыре разные точки. При $R=0$ результат состоит из центра.

\subsection*{Что показывает сглаживание и что измеряет программа}
При $x=4$ для того же отрезка Ву распределяет вклад между строками 2 и 3. Дробная часть $32/13-2=6/13$, поэтому нижнему пикселю достаётся вес $7/13$, верхнему --- $6/13$. Ближайший к идеальной прямой пиксель ярче выражен. Сумма весов равна единице.

При одинаковой длине два алгоритма могут иметь одинаковый порядок сложности и разное время. Ву создаёт до двух пикселей на шаг, а Брезенхем для отрезка --- один. В замерах приложения также учитывается создание объектов Python. Поэтому для сравнения важно назвать и длину отрезка, и то, какие операции вошли в замер.

\clearpage
\section*{Лабораторная 4. Какая часть фигуры попала в окно}
Отсечение находит геометрическое пересечение фигуры с окном. Граница включается в результат, поэтому при касании может остаться единственная точка.
\subsection*{Сазерленд --- Коэн. Отрезок и прямоугольник}
Каждый конец получает четыре бита. Слева направо они означают «выше», «ниже», «правее», «левее». Например, 0101 означает «ниже и левее». Для точки внутри окна или на его границе код равен 0000.
\begin{enumerate}
\item Побитовое ИЛИ кодов равно нулю: оба конца внутри. Прямоугольник выпуклый, поэтому принимаем весь текущий отрезок.
\item Побитовое И отлично от нуля: у концов есть общий внешний бит. Оба лежат за одной границей, и весь отрезок отбрасывается.
\item Иначе выбираем внешний конец, находим пересечение с прямой одной из указанных границ, заменяем конец и повторяем проверку.
\end{enumerate}
В варианте 13 используется явное уравнение $y=y_0+k(x-x_0)$, где $k=\Delta y/\Delta x$. Для вертикальной границы подставляем её $x$. Для горизонтальной находим $x=x_0+(y-y_0)/k$. У вертикального отрезка $x=x_0$. Горизонтальный отрезок вне допустимой полосы по $Y$ сразу отклоняется общим битом.

\textbf{Пример из программы.} Для $(-10,-4)\to(10,6)$ прямая имеет вид $y=1+x/2$. Окно ограничено $-6\le x\le6$, $-3\le y\le5$. Начальные коды 0101 и 1010. При $y=-3$ получаем $(-8,-3)$: это продолжение нижней стороны, точка всё ещё левее окна. Подстановка $x=-6$ даёт $(-6,-2)$. Второй конец проходит через $(8,5)$ и заменяется на $(6,4)$. Теперь оба кода нулевые. Для прямоугольника число проверок одного отрезка ограничено константой.
\subsection*{Сазерленд --- Ходжман. Выпуклый многоугольник}
Выпуклое окно есть пересечение внутренних полуплоскостей его рёбер. Обходим границы окна по очереди. На каждом проходе берём список вершин, оставшийся после предыдущего прохода. Проверяем каждое ребро $P\to Q$, включая переход от последней вершины к первой.
\begin{center}\begin{tabular}{lll}\toprule
Откуда & Куда & Что добавить в новый список\\\midrule
Изнутри & Внутрь & Конечную вершину $Q$\\
Изнутри & Наружу & Пересечение $I$\\
Снаружи & Внутрь & Пересечение $I$, затем $Q$\\
Снаружи & Наружу & Ничего\\\bottomrule\end{tabular}\end{center}
Для ребра окна $A\to B$ считаем $d(P)=\sigma[(B_x-A_x)(P_y-A_y)-(B_y-A_y)(P_x-A_x)]$. При обходе окна против часовой стрелки $\sigma=1$, по часовой $\sigma=-1$. Условие $d(P)\ge0$ означает, что точку сохраняем. При переходе через границу $t=d(P)/(d(P)-d(Q))$ и $I=P+t(Q-P)$. Формула следует из того, что вдоль ребра $d$ меняется линейно и в пересечении равна нулю.

\textbf{Почему вершин может стать больше?} Срез одного угла убирает одну вершину и добавляет две точки пересечения. В примере программы исходный четырёхугольник после четырёх границ даёт последовательно 5, 5, 6 и 6 вершин.

\textbf{Как читать экран.} Синий пунктир показывает исходную фигуру, оранжевый контур --- окно, зелёный --- видимую часть. Масштаб меняет только отображение. Формулы $X=o_x+sx$, $Y=o_y-sy$ переводят мировые координаты в экранные. Минус нужен потому, что экранная ось $Y$ направлена вниз.

\clearpage
\section*{Лабораторная 4. Как объяснить результат на экране}
\subsection*{Что означают две группы битов в строке таблицы}
Слева от косой черты записан код первого конца, справа --- второго. Строка показывает состояние \textbf{до} действия. Для отрезка $(-10,-4)\to(10,6)$ из примера получаем:
\begin{center}\begin{tabular}{lll}\toprule
Коды концов & Действие & Как меняется отрезок\\\midrule
0101 / 1010 & Снизу & Первый конец переходит в $(-8,-3)$\\
0001 / 1010 & Слева & Первый конец переходит в $(-6,-2)$\\
0000 / 1010 & Сверху & Второй конец переходит в $(8,5)$\\
0000 / 0010 & Справа & Второй конец переходит в $(6,4)$\\
0000 / 0000 & Принят & Видимая часть $(-6,-2)\to(6,4)$\\\bottomrule\end{tabular}\end{center}
\textbf{Почему нулевого побитового И мало для принятия?} Возьмём в том же окне отрезок $(-8,5)\to(0,9)$. Первый конец левее окна, второй выше, общих битов нет. Но на левой границе $x=-6$ ордината уже равна 6, тогда как верх окна имеет $y=5$. После замены первого конца оба конца оказываются выше окна, и отрезок отклоняется.

\textbf{Почему из шести отрезков видны пять?} Отрезок $(-9,3)\to(-7,6)$ целиком слева от окна. У отрезка $(-10,0)\to(-6,4)$ виден ровно один конец. Зелёная точка на левой стороне окна обозначает это касание и учитывается в числе видимых отрезков.

\subsection*{Один проход по многоугольнику вручную}
Исходный контур имеет вершины $A=(-8,-2)$, $B=(0,-6)$, $C=(8,-2)$ и $D=(0,6)$. При отсечении по нижней стороне $y=-3$ сохраняются точки с $y\ge-3$. Удобно взять $d(P)=P_y+3$: положительный множитель в формуле векторного произведения сократится при вычислении $t$.
\begin{enumerate}
\item $D\to A$: обе вершины внутри. Записываем $A$.
\item $A\to B$: выходим наружу. Значения $d(A)=1$ и $d(B)=-3$ дают $t=1/4$. Записываем $A+(B-A)/4=(-6,-3)$.
\item $B\to C$: входим внутрь. Получаем $t=3/4$, записываем пересечение $(6,-3)$ и вершину $C$.
\item $C\to D$: обе вершины внутри. Записываем $D$.
\end{enumerate}
На выходе пять вершин. Две новые точки соединяются ребром вдоль $y=-3$. Этот список целиком передаётся следующей границе. После всех четырёх проходов остаётся шестиугольник площадью $56.5$: из площади окна $9\cdot7$ вычитаются два угловых треугольника площадей $2$ и $4.5$.

\textbf{Зачем требуется выпуклость окна?} У выпуклого окна каждое ребро задаёт полуплоскость, содержащую всё окно. Для окна с вогнутым углом продолжение стороны может отрезать часть самого окна. Именно выпуклость обосновывает последовательную обработку границ.

\end{document}
